% \renewcommand{\arraystretch}{1.35}
\setlength{\tabcolsep}{8pt}

\begin{table}
\centering
\begin{tabular}{
    >{\raggedright\arraybackslash}m{0.15\textwidth}
    c
    >{\raggedright\arraybackslash}m{0.3\textwidth}
}
\toprule
\centering Name
&
\centering Rewrite rule
&
\centering Description
\tabularnewline
\midrule
    \zxlabel{cc}{Spider fusion ($f$)} &
    \begin{ZX}
        \leftManyDots{n_1} \zxZ{\alpha} \ar[dr,(=15] \ar[dr,)=15] \ar[dr,3 dots] \rightManyDots{m_1} &  \\
        & \leftManyDots{n_2} \zxZ{\beta} \rightManyDots{m_2}
    \end{ZX} = \begin{ZX}
        \leftManyDots{\begin{array}{c}n_1\\+\\n_2\end{array}} \zxZ{\alpha + \beta} \rightManyDots{\begin{array}{c}m_1\\+\\m_2\end{array}}
    \end{ZX}&
    Spiders of the same color connected by any number of wires fuse and their phases add.
\tabularnewline
\midrule
    \zxlabel{id}{Identity removal ($id$)} &
    \begin{ZX}
        \rar & \zxZ{} \rar &
    \end{ZX} = \begin{ZX}
        \rar & \zxN{} \rar & \zxN{}
    \end{ZX} &
    Spiders with one input, one output, and phase 0 are the identity matrix and can be removed.
\tabularnewline
\midrule
    \zxlabel{hh}{Hadamard cancellation ($hh$)} &
    \begin{ZX}
        \rar & \zxH{} \rar & \zxH{} \rar &
    \end{ZX} = \begin{ZX}
        \rar & \zxN{} \rar &
    \end{ZX} &
    Two connected Hadamards cancel each other.
\tabularnewline
\midrule
    \zxlabel{$\pi$}{$\pi$-commutation ($\pi$)} &
    \begin{ZX}
        \zxN{} & \zxN{} & \zxN{} & \zxN{} & \\
        \zxN{} \rar & \zxX{\pi} \rar & \zxZ{\alpha}
            \ar[ur,)]
            \ar[dr,(]
            & \zxN{} \ar[3 vdotsr={$n$}] & \\
        \zxN{} & \zxN{} & \zxN{} & \zxN{} &
    \end{ZX}
    = $e^{i\alpha}$
    \begin{ZX}
        \zxN{} & \zxN{} & \zxX{\pi} \rar & \\
        \zxN{} \rar & \zxZ{-\alpha}
            \ar[ur,)]
            \ar[dr,(]
            & \zxN{} \ar[3 vdotsr={$n$}] & \\
        \zxN{} & \zxN{} & \zxX{\pi} \rar & 
    \end{ZX}
    &
    A $\pi$ phase copies through a spider of the opposite color and flips its phase.
\tabularnewline
\midrule
    \zxlabel{cp}{State copy ($cp$)} &
    \begin{ZX}
        \zxN{} & \zxN{} & \zxN{} & \\
        \zxX{a\pi} \rar & \zxZ{\alpha} \ar[ur,)] \ar[dr,(] & \\
        \zxN{} & \zxN{} & \zxN{} & 
    \end{ZX} = $e^{ia \alpha}  \sqrt{2}^{-(n-1)}$
    \begin{ZX}
        \zxX{a\pi} \rar & \\
        \zxN{} \ar[3 vdotsr={$n$}] & \\
        \zxX{a\pi} \rar &
    \end{ZX} &
    A $\pi$ phase copies through a spider of the opposite color and flips its phase.
\tabularnewline
\midrule
    \zxlabel{cc}{Color change ($cc$)} &
    \begin{ZX}
        \zxN{} & \zxN{} \ar[dr,H,)={15}] & \zxN{} & \zxN{} & \zxN{} \\
            \zxN{} 
            & \zxN{} \ar[3 vdots={$n$}] 
            & \zxZ{\alpha} 
              \ar[ur,H,)={15}] \ar[dr,H,(={15}] 
            & \zxN{} \ar[3 vdotsr={$m$}] 
            & \zxN{} \\
        \zxN{} & \zxN{} \ar[ur,H,(={15}] & \zxN{} & \zxN{} & \zxN{}
    \end{ZX} =
    \begin{ZX}
        \zxN{} & \zxN{} \ar[dr,)={15}] & \zxN{} & \zxN{} & \zxN{} \\
            \zxN{} 
            & \zxN{} \ar[3 vdots={$n$}] 
            & \zxX{\alpha} 
              \ar[ur,)={15}] \ar[dr,(={15}] 
            & \zxN{} \ar[3 vdotsr={$m$}] 
            & \zxN{} \\
        \zxN{} & \zxN{} \ar[ur,(={15}] & \zxN{} & \zxN{} & \zxN{}
    \end{ZX} &
    Commuting Hadamards through a spider changes the color.
\tabularnewline
\midrule
    \zxlabel{b}{Bi-algebra ($b$)} &
    \begin{ZX}
        \zxN{} & \zxN{} \ar[dr,)={15}] & \zxN{} & \zxN{} & \zxN{} & \zxN{} \\
            \zxN{} 
            & \zxN{} \ar[3 vdots={$n$}] 
            & \zxX{} \rar
            & \zxZ{}
              \ar[ur,)={15}] \ar[dr,(={15}] 
            & \zxN{} \ar[3 vdotsr={$m$}] 
            & \zxN{} \\
        \zxN{} & \zxN{} \ar[ur,(={15}] & \zxN{} & \zxN{} & \zxN{} & \zxN{}
    \end{ZX} = $\sqrt{2}^{(n-1)(m-1)}$
    \begin{ZX}
        & \zxN{} \rar & \zxZ{} \rar \ar[ddr] & \zxX{} \rar & \zxN{} & \\
        &\zxN{} \ar[3 vdots={$n$}] & \zxN{} & \zxN{} & \zxN{} \ar[3 vdotsr={$m$}] & \\
        & \zxN{} \rar & \zxZ{} \rar \ar[uur] & \zxX{} \rar & \zxN{} &
    \end{ZX} &
    An adjacent phase 0 X and Z spider can commute past one another to create a complete bipartite graph with sides of the opposite color.
\tabularnewline
\midrule
    \zxlabel{hf}{Hopf rule ($hf$)} & 
    \begin{ZX}
        \zxN{} \rar & \zxZ{} \ar[r,(] \ar[r,)] & \zxX{} \rar & \zxN{}
    \end{ZX} = $\frac{1}{2}$
    \begin{ZX}
        \zxN{} \rar & \zxZ{} & \zxX{} \rar & \zxN{}
    \end{ZX} &
    If phase 0 spiders of opposite color are connected by multiple wires, we may remove the wires pairwise.
\tabularnewline
\bottomrule
\end{tabular}
\caption{ZX-Calculus Rewrite Rules}
\label{table:zx-rules}
\end{table}